\documentclass{article}
\input{../../lib.sty}

\title{\texttt{fn ResolvedRepetitions::sample}}
\author{Michael Shoemate}
\date{}

\begin{document}
\maketitle

Proves soundness of \texttt{fn ResolvedRepetitions::sample} in \texttt{mod.rs}.

\section{Hoare Triple}
\subsection*{Precondition}
\texttt{resolved} satisfies the proof definition of \texttt{ResolvedRepetitions}.

\subsection*{Pseudocode}
\lstinputlisting[language=Python,firstline=2,escapechar=|]{./pseudocode/resolved_repetitions_sample.py}

\subsection*{Postcondition}
\begin{theorem}
\label{postcondition}
For any valid \texttt{resolved}, \texttt{resolved.sample()} either returns \texttt{Err(e)}
due to propagated cast or sampler errors, or returns an exact draw from the repetition law
specified by \texttt{resolved}.
\end{theorem}

\section{Proof}
The function has three cases.

\begin{itemize}
    \item If \texttt{resolved = Poisson \{ mean \}}, line \ref{line:poisson} invokes
    \rustdoc{traits/samplers/poisson/fn}{sample\_poisson}. By its postcondition, the return value
    is distributed as \texttt{Poisson(mean)}.
    \item If \texttt{resolved = NegativeBinomial \{ eta: 0, x \}}, then after the exact cast on
    line \ref{line:x-rational}, line \ref{line:logarithmic} invokes
    \rustdoc{traits/samplers/logarithmic/fn}{sample\_logarithmic\_exp}. By its postcondition, the
    return value is distributed as the logarithmic member of Definition 1.
    \item If \texttt{resolved = NegativeBinomial \{ eta, x \}} with \texttt{eta > 0}, then after
    the exact casts on lines \ref{line:x-rational} and \ref{line:eta-rational}, line
    \ref{line:negative-binomial} invokes
    \rustdoc{traits/samplers/negative_binomial/fn}{sample\_truncated\_negative\_binomial\_rational}.
    By its postcondition, the return value is distributed as the Definition 1
    truncated-negative-binomial law.
\end{itemize}

These are the only branches of the function, which establishes the postcondition.

\end{document}
